\section*{Introduction}

The aim of this lab is to analyse the power distribution network (\textbf{PDN}) of a real test case provided by STMicroelectronics, composed of
a board, a four-layer package and a die, and to use this analysis to design its decoupling. The switching activity of the die draws a noisy
current from the supply, and the impedance of the PDN converts this current into a voltage ripple at the die terminals.

\noindent\\ When a peak of the PDN impedance coincides with a strong component of the current spectrum, as happens here around $100\,\text{MHz}$, the
ripple becomes large and the correct operation of the device is at risk. Decoupling capacitors mounted on the package mitigate this resonance and
shift it to less critical frequencies.

\noindent\\ The lab is divided into the following activities:
\begin{enumerate}
  \item \textbf{assembling the complete PDN SPICE netlist} from the provided board, package and die models, with all the port references connected
  to the global ground;
  \item \textbf{computing, in frequency domain, the PDN impedance} seen from the die interface and identifying its resonances;
  \item \textbf{validating the SPICE models of the decoupling capacitors} provided by the vendor against the vendor frequency-domain data;
  \item \textbf{performing a transient analysis of the PDN} with the supplied current load and evaluating the voltage ripple without decoupling capacitors;
  \item \textbf{connecting the decoupling capacitors and finding the configuration that minimises the voltage ripple}.
\end{enumerate}

\noindent\\ The final goal is the last one. The package offers five locations (\textit{C01--C05}) for the capacitors, and each of them can be populated
with one of six available values, from $1.5\,\text{nF}$ to $10\,\text{nF}$, so that the possible configurations are $6^5 = 7776$. Simulating all of
them in the time domain is not practical, since a single transient simulation of the whole PDN over the supplied current waveform takes several minutes,
and the complete search would require days of computation. For this reason we exploited the linearity of the PDN. The network without capacitors is
characterised once, through six AC analyses, by its impedance matrix evaluated at the frequencies of the spectral lines of the current. The capacitors
are then added analytically, for each configuration, and the voltage ripple is reconstructed from the current spectrum. In this way seven AC
simulations replace the 7776 transient simulations, and the best configurations found are finally confirmed with a transient simulation.

\noindent\\ To implement this approach we used an AI assistant (Claude, developed by Anthropic) as a means of support. It helped us to create the
MATLAB and Python scripts used to post-process the simulation data, to validate the capacitor models and to find automatically the optimal
configuration, the latter being based on the circuit-theory equations that are the core of the method (the impedance matrix of the multiport PDN
terminated by the capacitors). The circuit simulations were run by us in LTspice, and the results provided by the scripts were cross-checked against
these simulations.

\noindent\\ The report is organised in five sections, one for each of the activities listed above.